% BSDT UK Patent Application — Specification
% Title: Method and System for Detecting Hidden Risk Concentrations in
%        Financial Networks Using Tensor Decomposition
% Applicant: Odeyemi Olusegun Israel
% Date: March 2026

\documentclass[12pt,a4paper]{article}
\usepackage[margin=2.5cm]{geometry}
\usepackage{amsmath,amssymb}
\usepackage{graphicx}
\usepackage{enumitem}
\usepackage{hyperref}
\usepackage{titlesec}
\usepackage{fancyhdr}
\usepackage{booktabs}
\setlength{\headheight}{14.5pt}

\pagestyle{fancy}
\fancyhf{}
\rhead{Patent Application}
\lhead{BSDT --- Odeyemi O.I.}
\rfoot{Page \thepage}

\titleformat{\section}{\large\bfseries}{\thesection.}{0.5em}{}
\titleformat{\subsection}{\normalsize\bfseries}{\thesubsection.}{0.5em}{}

\begin{document}

\begin{center}
{\LARGE\bfseries PATENT SPECIFICATION}\\[1em]
{\Large Method and System for Detecting Hidden Risk Concentrations\\
in Financial Networks Using Tensor Decomposition}\\[2em]
{\large Applicant: Odeyemi Olusegun Israel}\\[0.5em]
{\large Inventor: Odeyemi Olusegun Israel}\\[0.5em]
{\large Date of Filing: 4 March 2026}\\[0.5em]
{\large Application Number: [TO BE ASSIGNED]}\\[0.3em]
{\normalsize \textit{Filed pursuant to the Patents Act 1977}}
\end{center}

\vspace{2em}

% ============================================================
\section{TECHNICAL FIELD}
% ============================================================

The present invention relates to the field of machine learning and 
anomaly detection, and more particularly to methods and systems for 
diagnosing and correcting structural failure modes in deployed classification 
models, with application to fraud detection in financial, aerospace, and 
medical domains.

% ============================================================
\section{BACKGROUND OF THE INVENTION}
% ============================================================

\subsection{Technical Problem}

Machine learning models deployed for fraud detection inevitably fail to 
detect certain classes of fraudulent activity. These failures are structural 
in nature---they arise from specific, identifiable deficiencies in how the 
model represents the data space---rather than from random noise. In 
production environments, these failures manifest as false negatives: 
fraudulent transactions that the model classifies as legitimate.

Current approaches to addressing model failures include:

\begin{enumerate}[label=(\alph*)]
\item \textbf{Model retraining:} Collecting new labelled data and 
retraining the model from scratch. This is expensive, requires fresh 
labelled data (which may not be available), and may disrupt certified 
production systems that have undergone regulatory approval.

\item \textbf{Ensemble methods:} Combining multiple models to reduce 
variance. While effective at reducing random errors, ensembles do not 
diagnose \emph{why} specific instances are missed, and multiple models may 
share the same structural blind spot.

\item \textbf{Cost-sensitive reweighting:} Adjusting class weights or 
decision thresholds. This is a global adjustment that trades precision for 
recall uniformly, without distinguishing between different failure modes.

\item \textbf{Bias-variance decomposition:} The classical decomposition 
operates in \emph{error space} (aggregate model performance), not in 
\emph{feature space} (why a specific instance was missed). It provides no 
per-instance diagnostic information.
\end{enumerate}

None of these approaches provides a structured, per-instance diagnosis of 
\textit{why} a particular observation was missed, nor a targeted correction 
mechanism that addresses specific failure modes without modifying the 
underlying model.

\subsection{Prior Art Limitations}

Known prior art techniques for post-hoc model correction include conformal 
prediction (which provides calibrated prediction sets but does not diagnose 
failure modes), Bayesian uncertainty estimation (which quantifies overall 
prediction uncertainty but does not decompose it into actionable components), 
and SHAP/LIME-based explanation methods (which explain individual 
predictions but do not identify structural failure patterns or provide 
correction mechanisms).

% ============================================================
\section{SUMMARY OF THE INVENTION}
% ============================================================

The present invention provides a method and system for decomposing machine 
learning model failures into four near-orthogonal, measurable components 
that operate in feature space, and for applying targeted correction operators 
that recover missed detections without modifying the underlying model. The 
invention addresses the limitations of the prior art through the following 
technical contributions:

\begin{enumerate}[label=(\roman*)]
\item A four-component decomposition of model failure modes into 
Camouflage ($C$), Feature Gap ($G$), Activity Anomaly ($A$), and Temporal 
Novelty ($T$), each computed from the feature space of the observation and 
reference distributions.

\item A self-calibrating weighting mechanism that computes component weights 
using one or more of Mutual Information, Fisher Variance Ratio, Bayesian 
Online, or Gradient-based methods.

\item A composite blind-spot score that aggregates the four components into 
a single scalar indicator of model vulnerability.

\item Post-hoc correction operators---including additive correction, 
logistic regression, quadratic surface fitting, and exponential gating---that 
apply targeted adjustments to model predictions based on the decomposed 
failure profile.

\item Verification of near-orthogonality of the four components by 
normalised mutual information (mean NMI $< 0.065$) and variance inflation 
factor (all VIF $< 1.45$), confirming that each component captures an 
independent aspect of model failure.
\end{enumerate}

% ============================================================
\section{DETAILED DESCRIPTION OF THE INVENTION}
% ============================================================

\subsection{Four-Component Decomposition}

The method operates on a frozen (unmodified) machine learning model $f$ that 
produces predicted probabilities $p(x) = f(x)$ for input observations $x$. 
The method computes four components for each observation, each measuring a 
distinct structural aspect of why the model may fail on that observation.

\subsubsection{Camouflage Component ($C$)}

The camouflage component measures how closely a potentially fraudulent 
observation resembles the legitimate data distribution. It is defined as:
\begin{equation}
C(x) = 1 - \frac{\|x - \mu_{\text{legit}}\|_2}{d_{\max}}
\end{equation}
where $\mu_{\text{legit}}$ is the mean of the legitimate class in the 
training data and $d_{\max}$ is the maximum distance from 
$\mu_{\text{legit}}$ observed in the training data. A value of $C(x)$ near 
1 indicates high camouflage---the observation closely mimics legitimate 
patterns, making it difficult for the model to distinguish from genuine 
transactions.

\subsubsection{Feature Gap Component ($G$)}

The feature gap component measures the proportion of features that carry 
negligible information for the observation:
\begin{equation}
G(x) = \frac{|\{j : |x_j| < \varepsilon\}|}{d}
\end{equation}
where $d$ is the total number of features and $\varepsilon = 10^{-8}$ is a 
near-zero threshold. A high value of $G(x)$ indicates that many features 
are uninformative for this observation, depriving the model of the 
discriminative signal it relies upon.

\subsubsection{Activity Anomaly Component ($A$)}

The activity anomaly component measures whether the transactional activity 
profile of the observation deviates from the activity profiles of previously 
detected fraud:
\begin{equation}
A(x) = \sigma\left(\frac{\log(1 + |\text{tx\_count}(x)|) - 
\mu_{\text{caught}}}{\sigma_{\text{caught}}}\right)
\end{equation}
where $\text{tx\_count}(x)$ is the transaction count associated with the 
observation, $\mu_{\text{caught}}$ and $\sigma_{\text{caught}}$ are the mean 
and standard deviation of the log-transformed transaction counts of 
previously detected (``caught'') fraudulent entities, and $\sigma$ is the 
logistic sigmoid function. A value near 1 indicates that the observation's 
activity level is anomalous relative to known fraud patterns.

This component has no direct analogue in the classical bias-variance 
decomposition. It captures a distinct axis of model vulnerability: fraud 
operating at activity scales the model has not learned to associate with 
fraud.

\subsubsection{Temporal Novelty Component ($T$)}

The temporal novelty component measures the degree to which the observation 
represents a pattern that is temporally novel relative to the training data:
\begin{equation}
T(x) = \sigma\left(\frac{1}{2}(\bar{m}(x) - 2)\right)
\end{equation}
where $\bar{m}(x)$ is the mean distance to the $k$ nearest temporal 
neighbours in the training set, normalised by the temporal span. A value 
near 1 indicates high temporal novelty---a pattern that emerged after 
training and that the model cannot be expected to recognise.

\subsection{Self-Calibrating Weight Computation}

The four component values for an observation are aggregated into a composite 
blind-spot score (designated MFLS --- Model Failure Likelihood Score) using 
weights that are computed automatically from the data:
\begin{equation}
\text{MFLS}(x) = \sum_{i \in \{C, G, A, T\}} w_i \cdot S_i(x)
\end{equation}

Four methods for computing the weights $w_i$ are provided:

\begin{enumerate}[label=(\alph*)]
\item \textbf{Mutual Information:} Weights are proportional to the mutual 
information between each component and the binary label (fraud vs. 
legitimate) on a held-out calibration set.

\item \textbf{Fisher Variance Ratio:} Weights are proportional to the 
ratio of between-class variance to within-class variance for each component.

\item \textbf{Bayesian Online:} Weights are updated sequentially using a 
Bayesian update rule as new labelled observations are received.

\item \textbf{Gradient-based:} Weights are optimised to maximise detection 
performance on a holdout set via gradient descent.
\end{enumerate}

The weights are normalised such that $\sum_i w_i = 1$.

\subsection{Correction Operators}

Given the four-component decomposition, the method applies correction 
operators to the frozen model's predicted probabilities to recover missed 
detections. Three correction operators are provided:

\subsubsection{Additive Correction}

The simplest correction adds a scaled composite score to the predicted 
probability for observations below a threshold:
\begin{equation}
p^*(x) = p(x) + \lambda \cdot \text{MFLS}(x) \cdot (1 - p(x)) \cdot 
\mathbb{1}[p(x) < \tau]
\end{equation}
where $\lambda$ is a correction strength parameter and $\tau$ is a threshold 
below which correction is applied.

\subsubsection{Logistic Regression Correction (SignedLR)}

A lightweight logistic regression is trained on the four BSDT components as 
input features, predicting the true label. The correction is:
\begin{equation}
p^*(x) = \text{clip}\left(p(x) + \beta_0 + \sum_{i} \beta_i S_i(x), 
\ 0, \ 1\right)
\end{equation}

\subsubsection{Quadratic Surface Correction (QuadSurf)}

A degree-2 polynomial ridge regression is fitted on the four components and 
their interactions, providing a richer correction surface:
\begin{equation}
p^*(x) = \text{clip}\left(p(x) + \phi_2([C, G, A, T])^\top 
\hat{\beta}, \ 0, \ 1\right)
\end{equation}
where $\phi_2$ denotes the degree-2 polynomial feature map that includes 
all terms $S_i$, $S_i^2$, and $S_i S_j$ for $i \neq j$. In the preferred 
embodiment, ridge regularisation with coefficient $\alpha = 1.0$ is applied.

\subsubsection{Exponential Gating Correction (ExpGate)}

A gated correction combines a quadratic surface with a sigmoid-activated 
gate:
\begin{equation}
p^*(x) = \text{clip}\left(p(x) + g(x) \cdot q(x), \ 0, \ 1\right)
\end{equation}
where $q(x)$ is the quadratic surface output and $g(x) = \sigma(v^\top 
[C, G, A, T] + b)$ is a learned sigmoid gate that modulates the correction 
strength.

\subsection{Seven Supplementary Features}

In an enhanced embodiment, seven supplementary features are derived from the 
core decomposition for use by the correction operators:
\begin{enumerate}
\item[$\varphi_1$--$\varphi_4$:] The four raw components $C$, $G$, $A$, $T$.
\item[$\varphi_5$:] The MFLS composite score.
\item[$\varphi_6$:] The predicted-vs-actual residual: $|p(x) - y|$ where 
$y$ is the true label (available only in supervised calibration).
\item[$\varphi_7$:] An interaction term: $C(x) \cdot A(x)$, capturing the 
joint effect of high camouflage with anomalous activity.
\end{enumerate}

\subsection{Orthogonality Verification}

The four components are verified to be near-orthogonal by three independent 
measures:
\begin{itemize}
\item \textbf{Normalised Mutual Information:} Mean NMI across all component 
pairs ranges from 0.051 to 0.064, well below the redundancy threshold of 
0.30.
\item \textbf{Variance Inflation Factor:} All VIF values are below 1.45, 
confirming negligible multicollinearity.
\item \textbf{Principal Component Analysis:} All four components carry 
meaningful variance; no component is dominated by others.
\end{itemize}

\subsection{Cross-Domain Applicability}

The method has been validated across seven datasets spanning four domains:
\begin{center}
\begin{tabular}{lll}
\toprule
\textbf{Domain} & \textbf{Dataset} & \textbf{Dominant Component} \\
\midrule
Blockchain & Elliptic & Temporal Novelty ($T$) \\
Blockchain & Ethereum Addresses & Temporal Novelty ($T$) \\
Banking & Credit Card Fraud & Activity Anomaly ($A$) \\
Banking & IEEE-CIS & Activity Anomaly ($A$) \\
Aerospace & Shuttle & Camouflage ($C$) \\
Medical & Mammography & Camouflage ($C$) \\
Pattern Recognition & Pendigits & Feature Gap ($G$) \\
\bottomrule
\end{tabular}
\end{center}

The domain-dependent variation in dominant components confirms that the 
decomposition captures genuine structural differences in failure modes 
across application domains.

\subsection{Performance Characteristics}

In the preferred embodiment, evaluated across 792 variant--model--dataset 
combinations:

\begin{center}
\begin{tabular}{ll}
\toprule
\textbf{Metric} & \textbf{Value} \\
\midrule
Best variant (QuadSurf) mean F1 & 0.787 ($+43\%$ relative improvement) \\
False positive reduction (QuadSurf) & 88.1\% (44,740 $\to$ 5,345) \\
Maximum recall recovery & $+84.4$ percentage points \\
Shuttle domain rescue: F1 & $0.134 \to 0.968$ \\
Shuttle domain FP reduction & 99.94\% \\
Zero-missed-fraud configurations & 10/12 model--dataset combos \\
Cross-domain credit card MFLS AUC & 0.885 \\
Cross-domain aerospace MFLS AUC & 0.982 \\
Cross-domain medical imaging MFLS AUC & 0.916 \\
\bottomrule
\end{tabular}
\end{center}

\subsection{Cross-Domain Transfer}

The damping mechanism derived from the BSDT framework transfers to physical 
systems. Specifically, the adaptive damping function:
\begin{equation}
\gamma^*(E) = \frac{E}{E + \theta}
\end{equation}
where $E$ is the blind-spot energy and $\theta$ is a damping threshold, has 
been validated as a viscosity regularisation term in Navier--Stokes fluid 
dynamics simulations, representing a novel finance-to-physics transfer 
direction.

% ============================================================
\section{CLAIMS}
% ============================================================

\noindent\textit{The following claims define the scope of protection
sought. Where a claim refers to another claim, the dependency is by claim
number. Features of any dependent claim may be combined with any
independent claim or with other dependent claims.}
\vspace{0.8em}

\begin{enumerate}[label=\textbf{\arabic*.}]

% --- Independent Method Claim ---
\item A computer-implemented method for detecting and correcting failure 
modes in a deployed machine learning classification model without modifying 
the model, the method comprising the steps of:
\begin{enumerate}[label=(\alph*)]
\item receiving a feature vector $x$ and a predicted probability $p(x)$ 
output by the deployed classification model;

\item computing a camouflage component $C(x)$ that measures the proximity 
of the feature vector to the centroid of a reference legitimate population 
in feature space;

\item computing a feature gap component $G(x)$ that measures the proportion 
of features in the feature vector having values below a near-zero threshold;

\item computing an activity anomaly component $A(x)$ that measures the 
deviation of a transactional activity metric of the observation from the 
activity profile of a reference population of previously detected anomalies;

\item computing a temporal novelty component $T(x)$ that measures the 
degree to which the observation is temporally distant from patterns present 
in the training data;

\item computing a composite blind-spot score by aggregating the four 
components using automatically calibrated weights; and

\item applying a correction operator to the predicted probability $p(x)$ 
based on values of one or more of the computed components to produce a 
corrected probability $p^*(x)$.
\end{enumerate}

% --- Independent System Claim ---
\item A system for detecting structural failure modes in a deployed machine 
learning classification model, the system comprising:
\begin{enumerate}[label=(\alph*)]
\item a decomposition module configured to compute, for each input 
observation, four near-orthogonal failure mode components comprising:
\begin{enumerate}[label=(\roman*)]
\item a camouflage component measuring proximity to a legitimate-class 
centroid;
\item a feature gap component measuring informational sparsity;
\item an activity anomaly component measuring deviation from detected-anomaly 
activity patterns; and
\item a temporal novelty component measuring temporal distance from training 
data patterns;
\end{enumerate}

\item a weighting module configured to compute component weights by one or 
more of mutual information, Fisher variance ratio, Bayesian online updating, 
or gradient-based optimisation;

\item an aggregation module configured to compute a composite blind-spot 
score from the weighted components; and

\item a correction module configured to apply one or more correction 
operators to predicted probabilities output by the deployed model, based on 
the computed component values, to produce corrected probabilities;
\end{enumerate}
wherein the system operates on a frozen deployed model without requiring 
modification or retraining of the deployed model.

% --- Dependent Claims ---
\item A method according to claim 1, wherein the camouflage component is 
computed as:
\[
C(x) = 1 - \frac{\|x - \mu_{\text{legit}}\|_2}{d_{\max}}
\]
where $\mu_{\text{legit}}$ is the centroid of the legitimate class and 
$d_{\max}$ is the maximum observed distance from the centroid.

\item A method according to claim 1, wherein the feature gap component is 
computed as:
\[
G(x) = \frac{|\{j : |x_j| < \varepsilon\}|}{d}
\]
where $d$ is the number of features and $\varepsilon$ is a near-zero 
threshold.

\item A method according to claim 1, wherein the activity anomaly component 
is computed as:
\[
A(x) = \sigma\left(\frac{\log(1 + |\text{activity}(x)|) - 
\mu_{\text{ref}}}{\sigma_{\text{ref}}}\right)
\]
where $\text{activity}(x)$ is a transactional activity metric, 
$\mu_{\text{ref}}$ and $\sigma_{\text{ref}}$ are statistics of a reference 
detected-anomaly population, and $\sigma$ is the logistic sigmoid function.

\item A method according to claim 1, wherein the correction operator is a 
quadratic surface correction comprising:
\[
p^*(x) = \text{clip}\left(p(x) + \phi_2([C, G, A, T])^\top 
\hat{\beta}, \ 0, \ 1\right)
\]
where $\phi_2$ is a degree-2 polynomial feature map of the four components 
and $\hat{\beta}$ are coefficients fitted by ridge regression.

\item A method according to claim 1, wherein the correction operator is an 
additive correction comprising:
\[
p^*(x) = p(x) + \lambda \cdot \text{MFLS}(x) \cdot (1 - p(x)) \cdot 
\mathbb{1}[p(x) < \tau]
\]
where $\lambda$ is a correction strength parameter and $\tau$ is a 
threshold.

\item A method according to claim 1, wherein the correction operator is an 
exponential gating correction comprising a quadratic surface output 
modulated by a learned sigmoid gate.

\item A method or system according to any preceding claim, wherein the 
four components are verified to be near-orthogonal by having normalised 
mutual information below 0.30 for all pairs and variance inflation factors 
below 2.0 for all components.

\item A method or system according to any preceding claim, wherein the 
dominant component varies by application domain, with temporal novelty 
dominating in blockchain domains, activity anomaly dominating in banking 
domains, and camouflage dominating in aerospace and medical domains.

\item A method or system according to any preceding claim, further 
comprising the step of computing seven supplementary features derived from 
the four components, comprising the four raw components, the composite 
score, a predicted-vs-actual residual, and an interaction term between the 
camouflage and activity anomaly components.

\item A method or system according to any preceding claim, wherein the 
adaptive damping function $\gamma^*(E) = E/(E + \theta)$ derived from the 
composite blind-spot energy is further applied as a viscosity regularisation 
term in a computational fluid dynamics simulation.

\end{enumerate}

% ============================================================
\section{DRAWINGS}
% ============================================================

\textit{[To be filed separately: Diagram of the four-component 
decomposition; flow diagram of the correction operator pipeline; 
orthogonality verification plots (NMI matrix, VIF chart); cross-domain 
performance comparison chart; schematic of the BSDT layer operating on a 
frozen model.]}

% ============================================================
\newpage
\section*{ABSTRACT}
% ============================================================

A computer-implemented method and system detect and correct structural
failure modes in deployed machine learning classification models without
model retraining. Model failures are decomposed into four near-orthogonal
components: a camouflage component measuring similarity of a fraudulent
observation to legitimate data; a feature gap component measuring
informational sparsity; an activity anomaly component measuring deviation
from established fraud detection patterns; and a temporal novelty component
measuring novelty relative to the training distribution. A composite score
aggregates the four components using self-calibrating weights. Post-hoc
correction operators apply targeted adjustments to predicted probabilities
based on the decomposed failure profile, recovering missed detections at
improvements of up to 84.4 percentage points of recall. The method operates
on any classification model without modifying its parameters, provides
interpretable attribution of failure cause, and is validated across
financial, aerospace and medical datasets.

\end{document}
